\section{Empirical Application}

\subsection{Empirical Objective and Research Questions}\label{sec:empirical_objective}
The empirical application is a predictive validation exercise. We ask whether a choice policy calibrated only to country-level preference anchors can recover meaningful cross-country variation on a separate survey instrument. The exercise does not estimate a human response distribution, identify causal effects, or recover individual preferences. Its purpose is to test whether information contained in aggregate anchors survives synthetic data generation and model training well enough to transport across instruments.

Three research questions organize the analysis. First, do the anchored policies encode the intended preference directions when evaluated on new choice scenarios? Second, do those directions correspond to independently observed cross-country patterns in the World Values Survey (\WVS)? Third, how does weight-based anchoring compare with simpler alternatives that use the same country information, including explicit persona prompts and mechanical transformations of the anchors? Together, these questions distinguish successful construction from external predictive value. A policy may faithfully encode its synthetic training signal while failing to transport to a human survey benchmark.

\subsection{Anchor Selection}\label{sec:anchor_selection}
A candidate anchor source must meet four conditions. It must be publicly documented before the application, provide comparable population-level measurements across countries, correspond to constructs that can be expressed through choices, and remain independent of the survey used for evaluation. The final condition rules out calibrating a policy to \WVS{} country means and then evaluating it on the same \WVS{} items. Broad country coverage is also important because the relevant evidence comes from cross-country variation rather than from the fit of a single national profile.

We selected the Global Preferences Survey (\GPS) because it best satisfies these conditions for the present application \cite{falk2018}. The \GPS{} reports country-level measures of trust, patience, risk taking, altruism, positive reciprocity, and negative reciprocity, built from survey items that were validated against incentivized experiments and collected in representative samples across 76 countries. Its country scores are available independently of \WVS{} Wave 7 and overlap with sixteen countries in our evaluation frame. Other sources could anchor future applications, including Schwartz values \cite{schwartz1992}, Hofstede dimensions \cite{hofstede2010}, Inglehart--Welzel cultural-map coordinates \cite{inglehart2005}, and the GLOBE study \cite{globe2004}. These alternatives measure broader cultural values or organizational practices, whereas the \GPS{} is directly organized around economic preference traits. We therefore use the \GPS{} as the anchor source and reserve the \WVS{} as an external benchmark.

\subsection{Application Design and Headline Findings}\label{sec:application_design}
The application covers the sixteen countries observed in both data sources. For each country, we train one adapter from its six-dimensional \GPS{} profile and evaluate the resulting choice policy on a theory-led battery of thirty single-choice \WVS{} items. Trust and patience provide the main cross-survey tests because their conceptual links can be stated in advance. The remaining preference dimensions are used to examine whether the training signal leaves a detectable imprint, with weaker claims about direct construct correspondence.

We compare each anchored policy with the unmodified base model, a persona-prompted version of the base model that names the country, and two mechanical mappings of the same anchors. The first mechanical benchmark converts each standardized anchor into the least-committal option distribution consistent with its implied mean; the second assigns the response that follows the sign of the anchor. These comparisons ask whether DPO and QLoRA contribute predictive structure beyond the information already present in the country scores. A restricted permutation test provides a separate check by reassigning the observed country anchors when computing the primary cross-country statistic. This comparison evaluates whether the realized country-to-anchor assignment carries information beyond generic synthetic tuning.

The experiments reveal two distinct forms of alignment. Persona prompting more closely reproduces the overall shape of the human response distributions, but it does not preserve the cross-country ordering of trust. The anchored policies recover that ordering and preserve the expected development gradient for patience, even though their response distributions remain farther from the human distributions on average. Additional tests show that prompting the trained adapters does not overturn the directional trust result, while temperature scaling mainly changes dispersion. These findings motivate the distinction used throughout the results section: distributional resemblance and directional transport are separate empirical properties.

This application should therefore be read as a test of whether declared aggregate moments can be converted into a reusable scoring instrument. We intend the resulting policies to complement human survey research by providing evidence about cross-country preference directions on new item batteries. They do not reconstruct within-country heterogeneity or establish that the adapters reproduce national cultures. Training-side replications with alternative synthetic banks, labeling rules, preference weights, and random seeds remain necessary to determine how sensitive the conclusions are to upstream design choices.
